\item Cuando partículas cargadas de alta energía se mueven a través de un medio transparente con una rapidez mayor que la rapidez de la luz en dicho medio, se produce una onda de choque, u onda de proa, de luz. Este fenómeno se llama efecto Cerenkov. Cuando un reactor nuclear se blinda mediante una gran alberca de agua, la radiación Cerenkov puede verse como un brillo azul en la vecindad del núcleo del reactor debido a electrones de alta rapidez moviéndose a través del agua (figura 16.39). En un caso particular, la radiación Cerenkov produce un frente de onda con un semiángulo de vértice de $53.0^\circ$. Calcule la rapidez de los electrones en el agua. La rapidez de la luz en el agua es $2.25 \times 10^8\text{ m/s}$.